% TOP_PROBLEM: {"id": 1401, "title": "Total Coloring Conjecture", "category_id": 3, "category": {"id": 3, "name": "graph_theory", "display_name": "Graph Theory", "description": "Problems involving graphs, networks, and their properties.", "slug": "graph-theory", "order_index": 3, "created_at": "2026-07-31T15:26:25.671Z"}, "status": "open"}
% FIRST_POSTED: null
% ENTRY_KIND: statement_only
% Imported from: batch07.tex; UnsolvedMath ID 1401
\documentclass[11pt]{article}
\usepackage[margin=25mm]{geometry}
\usepackage{fontspec}
\setmainfont{Latin Modern Roman}
\usepackage{amsmath,amssymb,mathtools}
\usepackage{unicode-math}
\setmathfont{Latin Modern Math}
\usepackage{xurl,hyperref}
\hypersetup{colorlinks=true,urlcolor=blue}
\setcounter{secnumdepth}{0}
\setlength{\parindent}{0pt}
\setlength{\parskip}{5pt}
\setlength{\emergencystretch}{3em}
\newcommand{\sref}[2]{\hyperlink{source:#1:#2}{[S#2]}}
\newcommand{\cataloguescope}[1]{\paragraph{Recorded target.}#1}
\begin{document}
% UnsolvedMath ID 1401; source catalogue code GRAPH-014
\setcounter{section}{1400}
\section{Total Coloring Conjecture}\label{1401}
{\small\sffamily UnsolvedMath ID 1401 · Graph Theory}
\subsection{Definitions and mathematical statement}

\paragraph{Shared notation.}$\mathbb N=\{1,2,\ldots\}$, and $\mathbb Z,\mathbb Q,\mathbb R,\mathbb C$ denote the integers, rationals, reals and complex numbers. Any other conventions are specified locally.


A graph $G=(V,E)$ is finite, simple and undirected: an edge is an unordered pair of distinct vertices. Its maximum degree is $\Delta(G)=\max_{v\in V}|\{e\in E:v\in e\}|$, with $\Delta=0$ for the empty graph. A total $k$-coloring is a function $c:V\sqcup E\to\{1,\ldots,k\}$ such that adjacent vertices receive different colors, edges sharing an endpoint receive different colors, and each edge receives a color different from each endpoint. The total chromatic number $\chi''(G)$ is the least number of colors allowing such a coloring; for the empty graph it is zero.
\paragraph{Statement recorded in the supplied source.}
Does every finite simple graph satisfy
\[
 \chi''(G)\le \Delta(G)+2?
\]
The assertion concerns all graphs, without planar, degree or forbidden-subgraph restrictions. \sref{1401}{2}
\subsection{Short English statement}
Can the vertices and edges of every finite simple graph be colored with at most two more colors than its maximum degree, while all adjacent or incident elements get different colors?
\subsection{Sources}
\begin{description}
\item[\hypertarget{source:1401:1}{S1}] ULAM.AI and the UnsolvedMath Contributors, UnsolvedMath: A Curated Collection of Open Mathematics Problems, record 1401 (GRAPH-014). CC BY 4.0. \url{https://huggingface.co/datasets/ulamai/UnsolvedMath}.
\item[\hypertarget{source:1401:2}{S2}] Rongjin Su, Gang Fang and Enqiang Zhu, The Total Coloring Conjecture Holds for Planar Graphs Without Three Special Subgraphs, arXiv:2507.12737v1 (2025), Section 1, Conjecture 1.1, and Section 2 simple-graph convention. \url{https://arxiv.org/pdf/2507.12737v1}.
\end{description}
\label{end:1401}
\end{document}
